\documentclass[10pt,aspectratio=169,fontset=none]{ctexbeamer}
\usetheme{Madrid}
\usecolortheme{seahorse}
\setCJKmainfont[Path=/Users/liyufeng/Library/Fonts/]{STHeiti-Alias.ttf}
\setCJKsansfont[Path=/Users/liyufeng/Library/Fonts/]{STHeiti-Alias.ttf}
\setCJKmonofont[Path=/Users/liyufeng/Library/Fonts/]{STHeiti-Alias.ttf}
\usepackage{booktabs}
\usepackage{amsmath}
\setbeamertemplate{navigation symbols}{}
\definecolor{acc}{HTML}{2A6FB0}
\definecolor{warn}{HTML}{C77D1E}
\setbeamercolor{frametitle}{bg=acc,fg=white}

\title{Schrödinger Bridge $\times$ RL $\times$ 机器人策略学习}
\subtitle{awesome\_study 全库综述：51 篇精读 + 2026 趋势雷达}
\author{awesome\_study}
\date{2026-09}

\begin{document}
\maketitle

\begin{frame}{主线问题：生成式策略的 $\log\pi$ 障碍}
\begin{columns}[T]
\column{0.48\textwidth}
\textbf{为什么生成式策略}
\begin{itemize}
  \item 机器人动作天然多模态，高斯 actor 把模式平均成无效动作
  \item Diffusion Policy（RSS 2023）：action chunk 条件去噪，12 任务平均 +46.9\%
  \item $\pi_0$ / RDT-1B 等 VLA 沿用 diffusion/flow 动作头
\end{itemize}
\column{0.48\textwidth}
\textbf{为什么 RL 卡住}
\begin{itemize}
  \item 策略梯度、SAC 熵项、PPO ratio 全要 $\log\pi(a|s)$
  \item 生成式策略动作是 SDE 终端边际：似然只有 ELBO 或需解 ODE
  \item 表达力与可优化性的结构性冲突
\end{itemize}
\end{columns}
\vspace{6pt}
\textbf{五条解法路线（2026 全部成型）}：逐步分解（DPPO/ReinFlow）· 路径空间（FLAC/GSB-MDPO）· 噪声空间（DSRL/LP-DS）· 条件化监督（RECAP/AWR）· 生成-选择（MVP/FMQ）
\end{frame}

\begin{frame}{仓库地图：七类 51 篇}
\small
\begin{tabular}{@{}p{0.24\textwidth}p{0.70\textwidth}@{}}
\toprule
\textbf{类别} & \textbf{篇目与定位} \\
\midrule
01 生成模型基础 (7) & DDPM · Score-SDE · FM · OT-CFM · Rectified Flow · MeanFlow · Diffusion Policy \\
02 SB 理论 (4) & Léonard · DSB · GSBM · SB Foundations —— 数学正典 $\to$ 神经化 $\to$ 广义化 $\to$ 教科书 \\
03 Bridge 算法 (7) & I$^2$SB · DSBM · SF$^2$M · SB Flow · ASBM · Soft-SB · PRISM \\
04 RL 基础 (5) & PPO · SAC · MDPO · GRPO · AWR —— SB$\times$RL 的接口层 \\
05 生成式策略 $\times$ RL (15) & DT · Diffuser · Diffusion-QL · FQL · CP · DPPO · Flow-GRPO · ReinFlow · DSRL · MP1 · DMPO · OFP · DBPO · OMP · FMQ \\
06 SB $\times$ RL 前沿 (4) & \textbf{FLAC · GSB-MDPO · RSBM · BDGxRL —— 选题主战场} \\
07 2026 雷达 (9) & MVP · MF-VLA · MFPO · ReactVLA · UCA-Flow · RECAP · LP-DS · DF-ExpEnse · BridgePolicy \\
\bottomrule
\end{tabular}
\vspace{4pt}

每篇配中文详细解读；19 篇前沿另配保版式中译 PDF。
\end{frame}

\begin{frame}{线索一：一步化的终局}
\textbf{步数演化}：DDPM(1000) $\to$ DDIM/FM（多步积分） $\to$ OT-CFM（少步依任务） $\to$ 蒸馏/Consistency Policy（1 或 3，需教师） $\to$ \textbf{MeanFlow 原生 1-NFE（免蒸馏）}
\vspace{4pt}

\textbf{关键机制}：学区间平均速度 $u(z_t,r,t)=\frac{1}{t-r}\int_r^t v\,d\tau$，恒等式
$u = v-(t-r)\tfrac{d}{dt}u$ 给出免教师回归目标；ImageNet 256 一步 FID 3.43。
\vspace{4pt}

\textbf{进入机器人（四篇代表）}：MP1（1-NFE 操纵，6.8ms）$\to$ DMPO（+PPO，104.2Hz 真机端到端）$\to$ OFP（免 JVP 自蒸馏，一步超 $\pi_{0.5}$ 十步）$\to$ DBPO（drift 固定点，105.2Hz 双臂）
\vspace{4pt}

\alert{判断：「多步 vs 一步」争论已结束；SB 必须卖一步化没有的东西。}
\end{frame}

\begin{frame}{线索二：路径空间——SB 给 RL 的两个恒等式}
\textbf{Girsanov 动能恒等式}
\[
\mathrm{KL}(\mathbb{P}^u\,\|\,\mathbb{Q}) = \mathbb{E}\Big[\int_0^1 \tfrac{1}{2\sigma^2}\|u\|^2\,dt\Big]
\qquad\text{同初态、同非零扩散及 Girsanov 条件下成立}
\]
FLAC：显式熵计算 $\to$ 动能预算 + 自动调节（ODE 仅有几何约束）。
\vspace{4pt}

\textbf{数据处理不等式}
\[
\mathrm{KL}(\mathbb{P}\,\|\,\mathbb{Q}) \ge \mathrm{KL}(\mathbb{P}_1\,\|\,\mathbb{Q}_1)
\qquad\text{压路径就压住了执行动作分布}
\]
GSB-MDPO：MDPO 邻近项搬进路径空间，加权 drift-MSE 一阶近似路径 KL。
\vspace{4pt}

\alert{格局：理论位已被占满；两篇以状态仿真为主——视觉操纵/VLA/真机全空。MFPO（ICML 2026）走似然近似路线正面叫板。}
\end{frame}

\begin{frame}{线索三：桥式先验——informative source 的红利}
\begin{itemize}
  \item \textbf{图像（I$^2$SB）}：退化图当边界而非条件，四类修复、九种设置有指标取舍，特定补全 2--10 NFE 接近最佳。起点带信息 $\to$ 路径短 $\to$ 少步稳。
  \item \textbf{导航（RSBM）}：学习先验 + $\epsilon$ 旋钮（随机条件桥 $\leftrightarrow$ 确定插值），室内仿真 3 步 92\% 成功率免蒸馏。$\epsilon$ 谱系 = SB 独有的 coverage--straightness 权衡。
  \item \textbf{操纵（BridgePolicy, ICML 2026）}：观测嵌入 SDE 从观测先验起步，52 仿真 + 5 真机任务均值领先对应基线，10 NFE（FlowPolicy 为 1）——顶会正面验证桥式起点（但只做了 BC 半场）。
  \item \textbf{跨域（BDGxRL)}：DSB 在线翻译源域转移 + 奖励调制，MuJoCo 跨域 18 个设置均值均领先 DARC 等基线；像素级 sim-to-real SB 至今无人做。
\end{itemize}
\end{frame}

\begin{frame}{主战场：SB$\times$RL 四条进路对照}
\footnotesize
\begin{tabular}{@{}p{0.15\textwidth}p{0.19\textwidth}p{0.19\textwidth}p{0.19\textwidth}p{0.19\textwidth}@{}}
\toprule
 & \textbf{FLAC} & \textbf{GSB-MDPO} & \textbf{RSBM} & \textbf{BDGxRL} \\
\midrule
问题定义 & GSB 动能正则替代熵计算 & 邻近更新进路径测度 & few-step 桥矫正 & 动力学差距 = unpaired 翻译 \\
$\log\pi$ 处理 & SDE 约束 KL；ODE 约束 $W_2$ & path-KL 作邻近项 & 不涉及（纯监督） & 不涉及（翻译后 SAC） \\
实验面 & DMC + HumanoidBench & 8 Playground + 6 Gym-MuJoCo & 室内仿真 3 步 92\% & MuJoCo 跨域 \\
最大短板 & 无操纵/真机 & 路径比率高方差 & 无 RL & 低维状态 \\
启示 & informative reference 增量 & 装进操纵 = 占两格 & $\epsilon$ 谱是核心资产 & 像素级窗口开着 \\
\bottomrule
\end{tabular}
\end{frame}

\begin{frame}{2026 雷达（2025-11 至 2026-08，9 篇）}
\textbf{少步族扩张}：MVP（2602.13810，边界约束）· MF-VLA（2603.01469）· ReactVLA（2606.14255，仿真 2 步/真机 5 步）· \alert{MFPO（2604.14698, ICML：两步 MeanFlow+max-ent，FLAC 正面对手）} · UCA-Flow（2608.16153，对 MP1 报称 +9.3pp，均值口径不符）
\vspace{6pt}

\textbf{RL 微调路线}：$\pi^*_{0.6}$/RECAP（2511.14759：优势条件化免策略梯度，困难叠衣/浓缩咖啡成功吞吐超两倍）· LP-DS（2606.01151, ICML：DSRL 修正 + \textbf{动作熵与模式覆盖评测}）· DF-ExpEnse（2606.19656）· BridgePolicy（2512.07212, ICML）
\vspace{6pt}

\textbf{五条趋势}：一步化范式化 · $\log\pi$ 四路线互相对表（path-space 独缺真机）· informative source 上主会但只做 BC · mode coverage 从口号走向指标 · 工业界与学术界在「绕开似然」上合流
\end{frame}

\begin{frame}{入库补检带来的两个新判断}
\textbf{(1) 一步时代 path-space 还剩什么}：三个可比较的偏移约束——GSB-MDPO（路径 KL）· LP-DS（latent 偏移）· FMQ（平均速度场，局部线性 Q 解 $u^*=u^{\mathrm{ref}}+\eta\,\nabla_aQ/\|\nabla_aQ\|$）。同初态、同非零扩散下，路径 KL 对应相对漂移的加权 $L^2$，\textbf{与 FMQ 平均速度约束的一步等价性仍需证明}。
\vspace{6pt}

\textbf{(2) PRISM 把 SB 差异化钉死在 few-step 区间}：不可见性原理——线性高斯复原模型中，最优预测器与合规细化网格使正噪声参考在无限步恢复同一后验；共享网格及最优解唯一等条件下，最优噪声谱 $\propto$ 后验残差方差谱 $\to$ informative source 的候选规则，动作迁移待验证。
\vspace{6pt}

\textbf{共性病理（OMP / MVP）}：速度/漂移回归在低模长区\textbf{梯度饥饿}；MeanFlow 恒等式缺边界条件\textbf{解不唯一}。桥匹配是否同病需检查实际目标——one-step SB policy 必须内置方向对齐 + 边界约束。
\end{frame}

\begin{frame}{空白与机会：五个无人占位的格子}
\small
\begin{tabular}{@{}clp{0.42\textwidth}l@{}}
\toprule
\# & \textbf{组合} & \textbf{零件（全部已有文献）} & \textbf{窗口} \\
\midrule
1 & \alert{One-step SB policy} & MeanFlow 恒等式 $\times$ 桥边界 $\times$ $\epsilon$ 谱 & 6--12 月 \\
2 & Few-step SB $\times$ path-RL & GSB-MDPO/FLAC 内核 $\times$ RSBM 骨架 $\times$ 操纵 & 12 月内 \\
3 & Mode coverage 基准 & LP-DS 熵估计 $\times$ $\epsilon$ 谱做实验变量 & 先做先赢 \\
4 & 像素级 sim-to-real SB & SB Flow 规模化 $\times$ BDGxRL 框架 & 竞争最少 \\
5 & RL 加权耦合 & advantage 注入 IMF/OT-CFM 耦合 & 纯算法空格 \\
\bottomrule
\end{tabular}
\vspace{6pt}

\textbf{风险}：「SB 帽子」贬值加速；path-space 纯理论位已饱和；MeanFlow 族半年 6 篇——对表要指路线而非单篇。
\end{frame}

\begin{frame}{结论}
SB$\times$RL 不是成熟工具，是有清晰 niche 的研究空间。一步化与 path-space 的理论位在 2026 上半年被迅速占领，但三个交叉格子——\textbf{one-step SB policy、few-step$\times$path-space 的操纵落地、像素级 sim-to-real 桥}——每个零件都有文献而组合无人做。
\vspace{6pt}

SB 的持久差异化资产只有两个：\textbf{informative source（边界即条件）}与 \textbf{$\epsilon$ 谱系（coverage--straightness 同旋钮）}；不落在这两点上的 SB 包装工作都会贬值。
\vspace{10pt}

\textbf{产物}：51 篇中文详细解读（Codex GPT-6 逐篇二审）· 英文 PDF 51 篇 + 19 篇中译 · 速览 / 开放问题 / 对比表 / 阅读路线 / 术语表 · HTML PPT + 本 PDF。
\end{frame}

\end{document}
